\section{Preference Optimization and Alignment}

\subsection{Direct Preference Optimization}

Post-training alignment is essential to ensure that language models generate responses that are not only fluent but also aligned with human preferences, safety guidelines, and logical standards. Traditionally, Reinforcement Learning from Human Feedback (RLHF) has relied on a two-step process: first, training a reward model on human comparison data using the Bradley-Terry preference model, and second, optimizing the policy model using Proximal Policy Optimization (PPO) against the learned reward. However, this process is computationally unstable, sensitive to hyperparameters, and requires maintaining multiple active networks during training.

Direct Preference Optimization (DPO) bypasses the need for a separate reward model by establishing an analytical relationship between the reward function and the optimal policy \parencite{rafailov2024directpreferenceoptimizationlanguage}. Under the Bradley-Terry preference model, the probability that response $y_w$ is preferred over response $y_l$ given prompt $x$ is defined as:
\begin{equation}
    P(y_w \succ y_l \mid x) = \sigma(r(x, y_w) - r(x, y_l)),
\end{equation}
where $r(x, y)$ is the latent reward function. By expressing the reward function in terms of the policy $\pi_\theta$ and a reference policy $\pi_{\text{ref}}$:
\begin{equation}
    r(x, y) = \beta \log \frac{\pi_\theta(y \mid x)}{\pi_{\text{ref}}(y \mid x)},
\end{equation}
the DPO objective minimizes the negative log-likelihood of the preference data directly:
\begin{equation}
    \mathcal{L}_{\text{DPO}}(\theta; \pi_{\text{ref}}) = -\mathbb{E}_{(x, y_w, y_l) \sim \mathcal{D}} \left[ \log \sigma \left( \beta \log \frac{\pi_\theta(y_w \mid x)}{\pi_{\text{ref}}(y_w \mid x)} - \beta \log \frac{\pi_\theta(y_l \mid x)}{\pi_{\text{ref}}(y_l \mid x)} \right) \right],
\end{equation}
where $\sigma$ is the sigmoid function and $\beta$ is a hyperparameter that controls the strength of the KL divergence regularization. 

In the legal domain, DPO provides a powerful mechanism for aligning Small Language Models (SLMs). Instead of general human preferences, preference pairs in legal reasoning are designed to distinguish between logical correctness and error. By setting the chosen response $y_w$ as the correct, structurally valid legal explanation and the rejected response $y_l$ as a student rollout containing reasoning errors (such as logical leaps or statutory misinterpretations), DPO forces the model to maximize the probability of correct reasoning paths while actively penalizing typical failure modes \parencite{li2026legaldrilldiagnosisdrivensynthesislegal}.

\subsection{Large Language Model as a Judge Paradigm}

Evaluating generated text in specialized domains like law presents a significant challenge. Traditional automated evaluation metrics, such as BLEU and ROUGE, rely on exact n-gram matching against reference texts. While useful for machine translation or simple summarization, these metrics are inadequate for legal reasoning. A legal explanation can be semantically correct and logically sound while using entirely different terminology and phrasing than the reference text. Conversely, a generated response could have a high ROUGE score but contain a single negating word (e.g., omitting "not" or misstating an article number) that renders the entire legal argument incorrect.

To overcome these limitations, the "Large Language Model (LLM) as a Judge" paradigm has emerged. This approach leverages a highly capable, instruction-tuned LLM (such as GPT-4) to act as an automated evaluator. The judge model is provided with a prompt containing a structured evaluation rubric, the input case facts, the ground-truth reference, and the model's generated output. It is then instructed to evaluate the response based on multi-dimensional criteria, including factual groundedness, logical coherence, and adherence to specific structural tags (e.g., thinking and reasoning tag blocks).

In the context of distillation and preference data construction, the LLM Judge plays a diagnostic role. It is used to analyze the on-policy rollouts generated by the student model and classify them into diagnostic categories such as correct, partially correct, risky, or incorrect. This detailed classification enables the automatic filtering of high-quality preference pairs: rollouts diagnosed with logical leaps or false assumptions are categorized as rejected, while their corresponding ground-truth references are categorized as chosen. This automated loop enables large-scale, high-quality preference alignment without requiring manual annotation \parencite{li2026legaldrilldiagnosisdrivensynthesislegal}.
